% Template for ICASSP-2021 paper; to be used with:
%          spconf.sty  - ICASSP/ICIP LaTeX style file, and
%          IEEEbib.bst - IEEE bibliography style file.
% --------------------------------------------------------------------------
\documentclass{article}
\usepackage{spconf,amsmath,graphicx}
\usepackage{booktabs}

% Example definitions.
% --------------------
\def\x{{\mathbf x}}
\def\L{{\cal L}}

% Title.
% ------
\title{Robust Wind Turbine Blade Segmentation from RGB Images in the Wild}

%
% Single address.
% ---------------
\name{Raül~Pérez-Gonzalo$^{1,2}$, Andreas~Espersen$^{2}$, Antonio~Agudo$^{1}$\thanks{This work has been supported by the Innovation Fund Denmark under 2021 ID1044-0044A and by the project MoHuCo PID2020-120049RB-I00 funded by MCIN/AEI/10.13039/501100011033.}}
\address{$^{1}$Institut de Robòtica i Informàtica Industrial, CSIC-UPC, Spain\\$^{2}$Wind Power LAB, Copenhagen, Denmark}
%
% Two addresses (uncomment and modify for two-address case).
% ----------------------------------------------------------
%\twoauthors
%  {A. Author-one, B. Author-two\sthanks{Thanks to XYZ agency for funding.}}
%	{School A-B\\
%	Department A-B\\
%	Address A-B}
%  {C. Author-three, D. Author-four\sthanks{The fourth author performed the work
%	while at ...}}
%	{School C-D\\
%	Department C-D\\
%	Address C-D}
%
\begin{document}
%\ninept
%
\maketitle
%
\begin{abstract}
With the relentless growth of the wind industry, there is an imperious need to design automatic data-driven solutions for wind turbine maintenance. As structural health monitoring mainly relies on visual inspections, the first stage in any automatic solution is to identify the blade region on the image. Thus, we propose a novel segmentation algorithm that strengthens the U-Net results by a tailored loss, which pools the focal loss with a contiguity regularization term. To attain top performing results, a set of additional steps are proposed to ensure a reliable, generic, robust and efficient algorithm. First, we leverage our prior knowledge on the images by filling the holes enclosed by temporarily-classified blade pixels and by the image boundaries. Subsequently, the mislead classified pixels are successfully amended by training an on-the-fly random forest. Our algorithm demonstrates its effectiveness reaching a non-trivial 97.39\% of accuracy. % which precedes a new hole filling procedure
\end{abstract}
%
\begin{keywords}
Blade Segmentation, Wind Turbine Inspections, BU-Net, Hole Filling. %Blade Inspections, Hole Filling.
\end{keywords}
%

\vspace{-0.15cm}
\section{Introduction} 
\vspace{-0.15cm}

Wind energy has demonstrated to be an excellent alternative energy source, thanks to being completely renewable and environmentally sustainable~\cite{windreport}. The increase in installed capacity urges to adapt to more demanding challenges in terms of cost and logistics~\cite{30capacity}, with special focus on developing non-destructive inspection manners for health monitoring. Drone inspections have become a fast and reliable setup for wind turbine diagnosis~\cite{drone}, capturing around 400 high-resolution images per turbine to design a repair plan. All in all results in high-volume data that encourages the development of automated, scalable and cost-effective solutions. 

To successfully design a data-driven maintenance framework for wind turbines in operation, current researchers are tackling distinct image-based problems such as blade defect detection~\cite{aidrone,maskrcnn} or autonomous navigation~\cite{realtime, automateddrone}. However, current solutions have several limitations and are not being applied yet to the industry. Hence, these challenges must be simplified by means of distinct low-level vision techniques. To this aim, this paper proposes addressing segmentation to reduce the complexity of these learning tasks.

%Segmenter: Transformer for Semantic Segmentation
%Image Segmentation Using Deep Learning: A Survey
\begin{figure}[t!]
%\begin{minipage}[b]{0.94\linewidth}
  \centering
\resizebox{8.4 cm}{!} {
  \centerline{\includegraphics[width=8.5cm]{Figures/introd.png}}}
  \vspace{-0.3cm}
%\end{minipage}
\caption{\textbf{Wind turbine blade images in the wild} with distinct complexity due to the variation of the landscapes, the blade shape and size, and the presence of shadows or dirt.}
\label{fig:introd}
\vspace{-0.55cm}
\end{figure}

Current state-of-the-art segmentation methods rely on learning-based approaches, usually employing an encoder-decoder architecture~\cite{deeplab,sw,cvpr2022} and, recently, an attention mechanism on the encoder~\cite{bmvc,attention2}. With a properly tailored architecture, these models have proven to successfully obtain high precision results. Unfortunately, they do not provide a segmentation algorithm that generalizes for newly acquired images~\cite{generic}. A generic solution is essential to deal with a large variety of images (see Fig.~\ref{fig:introd}) obtained on regular wind turbine inspections under a different setup. 

A common technique to enhance the segmentation mapping is including a post-processing step. DeepLab~\cite{deeplab} utilized Conditional Random Fields to fine-tune the localized segmented objects. DeepLab V3+~\cite{deeplabv3+} is further strengthened through superpixel merging in~\cite{generic}. Another interesting approach is ResNeSt~\cite{resnest}, which last stage incorporates a channel-wise attention with multi-path representation. However, these methods do not leverage our prior information about blade images, such as that the blade traverses the image or that images of the same inspection share image properties.  

This work proposes a novel segmentation algorithm that identifies the blade region for very diverse images in appearance. First, a preliminary mask is obtained by a customized U-Net model~\cite{unet}, whose architecture still inspires current techniques~\cite{unetformer}. Then, this preliminary result is improved through a set of post-processing steps: an extended hole filling routine from~\cite{fill} followed by an unsupervised random forest~\cite{randomforest}; refined again through another hole filling step. The hole filling relies on properly identifying the blade image borders and the random forest on underfitting the preliminary masks from the same blade surface. As a result, we obtain an accurate, robust and efficient model for blade segmentation that can handle a wide variety of images in the wild.

%\section{Related Work}

The introduction of Fully Convolutional Neural Networks~\cite{fcn} meant a great advance for the segmentation field. The authors employed an encoder-decoder architecture that only contains convolutional layers. The encoder constitutes the contracting path and the decoder the expansive path, where both are built on repetitive blocks. A landmark framework came with the U-Net architecture~\cite{unet}, where they managed to capture high-level context and local patterns by using skip connections. The underlying idea was to use the same feature maps that are obtained in the contraction to expand the component feature vector. To this aim, the output of each encoder block was stacked to the corresponding input of the decoder block, ensuring that the information that is lost during the downsampling is recovered. 

Distinct variants and applications have come up thanks to the U-Net~\cite{unet}. In V-Net~\cite{vnet}, they presented a deeper model by replacing the stride component by another convolutional layer and a novel loss based on the Dice coefficient to deal with pixel-class imbalances. Attention U-Net~\cite{attentionunet} extended the architecture by an attention grid-based gate for efficient localization. U-Net++~\cite{unet++} redesigned the skip connections with nested convolutional blocks to adapt the intermediate encoder features before stacking them to the decoder block. RefineNet~\cite{refinenet} is a multi-scaled input network that propagates the gradients through the residual connections, enabling the model to obtain high-resolution accurate predictions. Despite the appearance of all these new segmentation techniques, the U-Net~\cite{unet} still demonstrates its capacity to outperform current state-of-the-art techniques in some tasks, such as in satellite image segmentation~\cite{satellite}.

%\cite{hrnet}
%\cite{pathnet}

Other segmentation techniques that also rely on the encoder-decoder framework include a post-processing step to enhance their predicted mapping. DeepLab family~\cite{deeplab} introduced dilated convolutions, capable of enlarging the receptive field of their convolutional without increasing the numbers of weights. In their first work, they utilized fully-connected Conditional Random Fields to fine-tune the localized segmented objects. On the other hand, its later version DeepLab V3~\cite{deeplabv3} achieved encouraging results by combining multiple scales of dilated convolutions through pyramid pooling. DeepLab V3~\cite{deeplabv3} is further strengthened through superpixel merging in~\cite{generic}. Another interesting approach for facilitating the information flow between distinct scales is BPNet~\cite{realtime}. Its architecture is based on a symmetric multi-scale network where the successive layers are continuously merged. Its last stage incorporates an unary attention mechanism for context aggregation. 


%CHALLENGES FOR MY PROBLEM (DATA) AND WHY CURRENT SOLUTIONS DO NOT WORK FOR OUR PROBLEM

%per aixo, proposem un metode de segmentacio per resoldre aquest problemes. la segmentacio es basa en bla bla. Els metodes actuals de segmentacio fallen amb dades noves \cite{}, no son robustos. Per aixo, proposem un metode adaptat, que aprofita el knowledge que tenim de les imatges de pales per tal.
%\section{Wind Turbine-aware Segmentation}
\section{Wind Turbine Blade Segmentation}
\label{section_methods}

Let $\mathbf{I}$ be a $H \times W \times 3$ color image of a wind turbine blade that was captured from an arbitrary point of view and under uncontrolled lighting conditions. Our goal is to recover a $H \times W$ binary mask $\hat{\mathbf{S}}$ to indicate the region in the image that belongs to the blade. In our formulation, the mask is represented by the binary entries $\hat s_{h,w}$, where $h=\{1,\ldots,H\}$ and $w=\{1,\ldots,W\}$ denote the pixel coordinates in the mask.

Given the corresponding ground-truth mask $\mathbf{S}$, we propose a supervised learning framework composed of four modules to learn the mapping $\mathbf{I}_{i,j}\rightarrow{\mathbf{S}_{i,j}}$ for a $i$-th blade and $j$-th  image: 1) an encoder-decoder network algorithm with a tailored loss, 2) a hole filling step, 3) a random forest block and, 4) a latter hole filling step to refine the solution. Without loss of generality, for every blade surface there are considered up to $J$ images for the $i$-th blade.

%it may be observed blades of very different sizes and shapes

%-----------------------------------------------------------------
\subsection{Blade U-Net (BU-Net)}

Given an input RGB image $\mathbf{I}$, our BU-Net aims at classifying each input pixel into blade or background, i.e., obtaining the output mask ${\hat{\mathbf{S}}}^{BU}$ by exploiting the labels in $\mathbf{S}$. This module follows a standard U-Net architecture~\cite{unet}, and its last layer includes a $1 \times 1$ convolution to map the feature channels to the desired number of classes. The mask ${\hat{\mathbf{S}}}^{BU}$ is obtained by combining four predictions via soft voting, where each one corresponds to a flipping rotation. Additionally, a regularized focal loss is employed to adapt the learning task to our particular dataset. The two terms are detailed below.

%-------------------------------------------------------------
The \textit{Focal Loss} \label{sec:focal-loss} optimizes our segmentation mapping~\cite{focal} by penalizing low-confidence predictions. Moreover, to avoid benefiting the majority class (in our case, the background) relative to the minority class, we apply weights to each class: 

\vspace{-0.2cm}
\begin{gather} 
\label{focal_loss}
    \mathcal L_{f}(\mathbf{S}, {\hat{\mathbf{S}}}^{BU}) = -\sum_{h,w} \left[ \alpha \left(1-\sigma(\hat s^{BU}_{h,w})\right)^\gamma s_{h,w}\log\sigma(\hat s^{BU}_{h,w})  \right. \nonumber \\ 
   \left.  + (1 - \alpha) \left(\sigma(\hat s^{BU}_{h,w}))\right)^\gamma (1-s_{h,w})\log(1-\sigma(\hat s^{BU}_{h,w}))    \right],
\end{gather}
where $\sigma(\cdot)$ denotes the sigmoid function, $\gamma$ controls the confidence penalty; and $\alpha$ sets the class balance.  

%the output is coded by a sigmoid function $sigma(\cdot)$ :
%\begin{align}
   % \mathcal L_{f}(y, \hat y) = -\sum_{i,j} \left[ \alpha \left(1-\sigma(\hat y_{i,j})\right)^\gamma y_{i,j}\log\sigma(\hat y_{i,j})  \right. \nonumber \\ 
  % \left.  + (1 - \alpha) \left(\sigma(\hat y_{i,j}))\right)^\gamma (1-y_{i,j})\log(1-\sigma(\hat y_{i,j}))    \right],
%\end{align}
%where $\gamma$ is a focusing weight to control the confidence; and $\alpha$ is a coefficient to set the class balance. 

%Let $\sigma$ denote the sigmoid function, $y$ is the desired output, $\hat y$ is the output of the model and $i, j$ are the coordinates of the image pixels. The Focal loss \cite{focal} weights each pixel by the confidence we have in the prediction of that pixel. Moreover, to avoid benefiting the majority class (in our case, the background) relative to the minority class we apply weights to each class:


%The optimal values obtained in \cite{focal} were $\gamma = 2$, which controls the confidence, and $\alpha=0.25$, which controls the class balance. Thus, we based our experiments around these values.

The \textit{Contiguity Loss}~\cite{cont} minimizes the jaggies in the blade boundaries, forcing to detect contiguous objects:
\begin{align}
    \mathcal L_{c}({\hat{\mathbf{S}}}^{BU}) = \frac{1}{HW} \Big[ \sum_{h,w} \left(   \sigma(\hat s^{BU}_{h+1,w})-\sigma(\hat s^{BU}_{h,w})  \right)^2 \nonumber \\
+ \sum_{h,w} \left(    \sigma(\hat s^{BU}_{h,w+1})-\sigma(\hat s^{BU}_{h,w})   \right)^2 \Big]^{1/2} .
\end{align}


The \textit{Total Loss} is defined as a linear combination of the two previous losses modulated by the scalar $\lambda$:
\begin{equation}
\label{total_loss}
\mathcal L = \mathcal L_{f} + \lambda\mathcal L_{c}\text{ } .
\end{equation}
 As the output of this block is the probability for the blade class, i.e., $\sigma(\hat{\mathbf{S}}^{BU})$, it is quantized by a threshold of 0.4.

%(\alpha = 0.25, \gamma = 2)

%\subsection{Hole Filling by Raül}
%After training the model, the segmented image is enhanced through a set of post-processing steps. Let $f$ denote the BU-Net, $\mathbf{x}_{ij}$ the $j$-th image of the blade $i$-th (in each inspection, there are $M_i$ images for each blade $i$) and $f(\mathbf{x}_{ij})$ the corresponding segmented image. To enhance the robustness of the BU-Net outcome, $f(\mathbf{x}_{ij})$ is obtained by combining four predictions via soft voting. These predictions correspond to the BU-Net output of $\mathbf{x}_{ij}$, where $\mathbf{x}_{ij}$ has been rotated by $k \dfrac{\pi}{2}$ for $k=\{0,1,2,3\}$. %ensure the rotation invariance 

%A blade mask is a connected component and it should be uniform. As the output of the network may present some irregularities on the blade shape,  we fill in the holes of the blade image regions using the classical operation presented in \cite{fill}. As this method is able to only fill in holes that are completely surrounded by blade pixels, we ensure that the blade image borders (pixels that lay in the image borders that correspond to blade) are identified as such, using our prior knowledge that they are continuously connected. 

%As we know that the blade traverses the image, from left-to-right (horizontal orientation) or top-to-bottom (vertical orientation), we first calculate the blade orientation. It is classified as vertical (horizontal) if the accumulated gradient along the x-axis is higher (lower) than along the y-axis; we compare the absolute sum of the horizontal and vertical Sobel on the BU-Net mask. Once the orientation is identified, we can distinguish where the blade image borders are. Therefore, we ensure those borders are continuously connected by replacing the background pixels within the blade area by blade pixels.


%------------------------------------------------------------------------
\subsection{Hole Filling}
\label{hole_filling}

%After training the model, the segmented image is enhanced through a set of post-processing steps. Let $f$ denote the BU-Net, $\mathbf{x}_{ij}$ the $j$-th image of the blade $i$-th (in each inspection, there are $M_i$ images for each blade $i$) and $f(\mathbf{x}_{ij})$ the corresponding segmented image. To enhance the robustness of the BU-Net outcome, $f(\mathbf{x}_{ij})$ is obtained by combining four predictions via soft voting. These predictions correspond to the BU-Net output of $\mathbf{x}_{ij}$, where $\mathbf{x}_{ij}$ has been rotated by $k \dfrac{\pi}{2}$ for $k=\{0,1,2,3\}$. %ensure the rotation invariance 

The estimated segmentation $\hat{\mathbf{S}}^{BU}$ still contains some non-realistic regions, including some holes. We propose a {\em hole filling} step to reduce these artifacts. Our approach is able to not only fill in the holes surrounded by blade pixels, but also when they are surrounded by the image borders. By exploiting a spatial prior, our approach relies on localizing first the blade pixels in the image boundaries, i.e., the blade-borders.


Assuming that the blade traverses the image, we first compute the blade orientation by means of an accumulated gradient. Basically, it is classified as vertical (horizontal) if the accumulated gradient along the $x$-axis is higher (lower) than along the $y$-axis. Once the orientation is identified, we know the blade-border locations, so we can enforce their continuity by setting their pixels as blade. Finally, we proceed with the standard hole filling algorithm~\cite{fill} to obtain the output $\hat{\mathbf{S}}^{H}$. Both vertical and horizontal orientation examples are displayed in Fig.~\ref{fig:holefilling}. The hole filling step is applied twice: 1) considering as input the estimation in our BU-Net and, 2) using as input the output of our random forest algorithm that is introduced below. To differ between both hole filling estimations, we will use the notation $\hat{\mathbf{S}}^{H1}$ and $\hat{\mathbf{S}}^{H2}$, respectively. 


%As we know that the blade traverses the image, from left-to-right (horizontal orientation) or top-to-bottom (vertical orientation), we first calculate the blade orientation. It is classified as vertical (horizontal) if the accumulated gradient along the x-axis is higher (lower) than along the y-axis; we compare the absolute sum of the horizontal and vertical Sobel on the BU-Net mask. Once the orientation is identified, we can distinguish where the blade image borders are. Therefore, we ensure those borders are continuously connected by replacing the background pixels within the blade area by blade pixels.



\begin{figure}[t!] %[htb]
\begin{minipage}[b]{1.0\linewidth}
  \centering
  \resizebox{8.5 cm}{!} {
  \centerline{\includegraphics[width=8.5cm]{Figures/holefilling.png}}}
  \vspace{-0.25cm}
\end{minipage}
\caption{\textbf{Hole filling representation.} First, the blade orientation is detected to identify the blade image borders (in orange). Then, holes that edge with the image borders are filled. }
\label{fig:holefilling}
\vspace{-0.3cm}
\end{figure}
%By inferring the orientation of the blade, our hole filling algorithm is capable to fill holes that edge with the image borders. Orange borders represent the identified blade image borders.


%--------------------------------------------------------------------
\vspace{-0.2cm}
\subsection{Unsupervised Random Forest}
\label{random_forest}


Finally, we propose a simple random forest to denoise the estimation and to fix uncommon errors of a single image. The inability of previous steps to capture global context is tackled by learning which RGB pixels are usually mapped as foreground, providing a fine-grained boundary precision. To this end, we learn the mapping $\{\mathbf{I}_{i,j}\}_{j \in J} \rightarrow \{\hat{\mathbf{S}}^{H1}_{i,j}\}_{j \in J}$ given the $i$-th blade surface. This model ensembles tree estimators by averaging their probabilistic predictions to set whether a single RGB pixel and its local neighborhood is foreground:

\vspace{-0.5cm}
\begin{equation}
\label{eq:rf} 
P^{RF}\big(\hat s_{h,w}^{RF} \big) = \frac{1}{|RF|} \sum_{t \in RF} P^t\big( \hat s_{h,w}^{t} | \hat s_{h,w}^{H1}, \text{neighbors}(\hat s_{h,w}^{H1}) \big)\text{ } , \nonumber
\vspace{-0.2cm}
\end{equation}
where $RF$ is the set of decision trees $t$. It is worth pointing out that this approach enables our algorithm to adapt to the image features of a specific $i$-th blade. Once the model is trained, it is applied to obtain the estimation $\hat{\mathbf{S}}^{RF}_{i,j}$.

Lastly, we again apply a hole filling (section~\ref{hole_filling}) step to refine the solution, obtaining $\hat{\mathbf{S}}^{H2}$ as our final estimation $\hat{\mathbf{S}}$. 








 

%\subsubsection{Data Augmentation.} \label{sec:data-augmentation}

%To allow the network learn the desired invariance and robustness properties, two strategies are applied to leverage the available annotated data (Fig. \ref{fig:data-aug}): 

%\begin{itemize}
 %   \item Flipping and rotating the image data by multiples of $\dfrac{\pi}{2}$.
    
  %  \item Rotating the input image by an angle sampled in $[-5^{\circ}, 5^{\circ}]$ following the uniform distribution. Subsequently, the transformed image is cropped in the created areas with no pixel data and resized to its original size, so there is no need to extrapolate over these areas.

%\end{itemize}

%\begin{figure}[H]
%\centering
%\includegraphics[width=4.8in]{Figures/data_augmentation.png}
%\caption{Data augmentation example used in segmentation. The same operations are applied to the corresponding mask.}
%\label{fig:data-aug}
%\end{figure}





%============================================================
%============================================================




\section{Experimental Results} \label{sec:results}

%In this section, we show the effectiveness of our segmentation approach through both quantitative and qualitative evaluations as well as showing a comparison with competing techniques, providing the results for every step in our algorithm.

In this section, we demonstrate the efficacy of our segmentation approach by conducting comprehensive quantitative and qualitative analysis. Additionally, we present a comparative assessment against competing techniques, offering detailed results for each step implemented within our algorithm.


%We now present the segmentation performance and computational cost of our proposed algorithm. To start with,  we examine the best hyperparameters reported in the BU-Net and random forest steps. The reliability of the algorithm is analyzed qualitatively and quantitatively, comparing each post-processing step. We also discuss whether our proposed algorithm can outperform the original U-Net. Finally, we cover the running time spent for each algorithm step.


%----------------------------------------------------------------------
\vspace{-0.5cm}
%\vspace{-0.3cm}
\subsection{Dataset} \label{sec:data} 
\vspace{-0.1cm}

We present a dataset that consists of high-resolution images taken by drone and ground-based equipment from distinct locations of the blade. The images provide a large variety of blade sizes, shapes, illumination conditions, as well as the background to be observed (Fig.~\ref{fig:introd}). Another background variation arises from on-shore and off-shore wind turbines. For this challenging dataset, we propose 1,712 images for training, 120 for validation and 200 for test. As a segmentation ground-truth is not provided, we manually annotated all the images. To make a fair analysis, the three sets are composed of different windfarms and inspection campaigns, ensuring their independence and that the test properly generalizes. In particular, we randomly selected 20 images of each windfarm to generate the validation and test sets. On the other hand, the training data was selected from a pool of different blade images, prioritizing the ones that are more challenging. These are the images in the root, the tip or the max-cord (where the blade section is widest). All color images and masks were min-max normalized. In addition to that, flipping and cropping strategies were applied to learn the desired invariance and robustness properties.



%The segmentation data consists in high-resolution images taken from distinct locations of the blade with their corresponding ground truth masks, which have been annotated manually. The dataset contains pictures taken by drone and ground-base equipment, which mainly diverge on its background. Other background variation arises from on-shore and off-shore wind turbines. The expected relation between background-blade pixels is 2:1, but in some cases the blade area could be very reduced. The data include 1712 training, 119 validation and 200 test images. The images from these three sets are from different windfarms and inspection campaigns, ensuring their independence and that the test properly generalizes.

 






%----------------------------------------------------------------
\vspace{-0.5cm}
%\vspace{-0.3cm}
\subsection{Implementation Details} \label{sec:training-seg}
\vspace{-0.1cm}

The original input images were resized to $1024 \times 1024$ pixels. We use Adam solver~\cite{adam} with an initial learning rate of $10^{-4}$ and a mini-batch of 1. A custom scheduler is employed to reduce the learning rate when there is no improvement on the validation loss after three epochs. Early stopping is performed. The model is trained using a NVIDIA GeForce RTX 3080 GPU. The weights for the focal loss in Eq.~\eqref{focal_loss} are set to $\gamma=2$ and $\alpha=0.25$, and the weight in Eq.~\eqref{total_loss} $\lambda=1$.



%----------------------------------------------------------------
\vspace{-0.5cm}
%\vspace{-0.3cm}
\subsection{Quantitative Evaluation}
\vspace{-0.1cm}

%The segmentation performance per image is analyzed in terms of the accuracy, recall and F1-score. 

%These metrics are defined pixel-wise as: $\textrm{accuracy} = \dfrac{TP + TN}{p}$, $\textrm{recall} =  \dfrac{TP}{TP+FN}$ and $\textrm{F1-score} = \dfrac{1}{\frac{1}{\text{precision}} + \frac{1}{\text{recall}}}$, where $TP$ indicates the number of True Positive pixels, $TN$ denotes the number of True Negative pixels, and $FN$ is the number of False Negative pixels. Moreover, precision is defined as $\frac{TP}{TP+FP}$ with $FP$ the number of False Positive pixels. 

The segmentation performance per image is analyzed in terms of accuracy, recall and F1-score. %The accuracy corresponds to the proportion of pixels of the predicted mask that match the ground truth mask, i.e., the ratio of correctly classified pixels. The recall is the ratio of blade pixels that are correctly classified by the algorithm. The F1-score is the harmonic mean of precision and recall, where the precision measures the accuracy of the blade predictions. The F1-score combines both precision and recall into a single metric, giving much more weight to low values.



%\begin{equation} \label{eq:accuracy}
 % \textrm{accuracy} = \dfrac{TP + TN}{p}, 
%\end{equation}


%\begin{equation} \label{eq:recall}
%  \textrm{recall} =  \dfrac{TP}{TP+FN}
%\end{equation}


%\begin{equation} \label{eq:f1}
%  \textrm{F1-score} = \dfrac{1}{\frac{1}{\text{precision}} + \frac{1}{\text{recall}}}
%\end{equation}

%where $TP = $ number of True Positive pixels, $TN = $ number of True Negative pixels, and $FN = $ number of False Negative pixels. Moreover, precision $= \frac{TP}{TP+FP}$ with $FP = $ number of False Positive pixels.




%the proportion of correctly classified blade pixels divided by the number of blade pixels in the ground truth mask.
%recall = true positive rate = true blade pixel rate


%The highest validation accuracy of the BU-Net is obtained using the original U-Net architecture, with a training objective $\mathcal L_{focal}(\alpha = 0.25, \gamma = 2) + \lambda  \cdot \mathcal L_{contiguity} $ where $\lambda = 1$, and a data augmentation strategy of small rotations followed by the corresponding zoom.


\textbf{Hyperparameter Search.} Our random-forest step (section~\ref{random_forest}) is optimized in terms of the number of split branches and the number of decision trees, obtaining the highest accuracy with five tree estimators that have a maximum of four split branches. Additionally, when the local pixel neighbors are included in the input, we enhance the random forest masks. However, increasing the number of input neighbors does not increase the segmentation accuracy, thus, we took as input the local pixel and its most immediate neighbors. 

\textbf{Test Set Results.} The performance over the test set is compared between each algorithm step. In particular, we report the accuracy, recall and F1-score in Table~\ref{tab:seg-accuracy}. 

\vspace{-1.5cm}
By only employing our BU-Net, we obtain a notable outcome of 92.28\% of accuracy. However, this result is not sufficiently good for blade assessments, because the recall of 85.73\% could compromise the detection of defects. For instance, if a large area of the blade is masked as background, it could cover the region where there are structural damages, which is a severe defect that cannot be overlooked.

By incorporating the first hole-filling step, we substantially improve performance. Namely, the recall increases to 92.08\%, which represents the highest gain of all the post-processing steps. This huge difference is due to the BU-Net being able to successfully capture the blade borders, but not the inner region of the blade. Further algorithm steps are necessary to yield further improvement, reaching in the end a novel performance of 97.39\% of accuracy and 93.35\% of recall. On balance, and as it can be seen in Table~\ref{tab:seg-accuracy}, every step is key to sort out the problem properly, obtaining a final accurate solution that cannot be achieved directly by our BU-Net. 

Fig.~\ref{fig:boxplot} analyzes the performance dispersion of each step. We observe the skewness of the BU-Net performance: rather spread along the tail with lower performance, containing many instances with poorer accuracy and recall compared with their median. As we include the different steps in our algorithm, that dispersion is reduced notably. Therefore, despite not substantially improving the mean performance, the subsequent steps are the cornerstone to ensure robustness.

\begin{table}[t!]
\centering
\caption{\textbf{Test results after each segmentation step.} }
    \vspace{-0cm}
    \label{tab:seg-accuracy}
%\begin{tabular}{p{1.95cm}p{1.0cm}p{1.2cm}p{1.6cm}p{1.1cm}}
\resizebox{8.5 cm}{!} {
\begin{tabular}{ccccc}
\toprule
     \multicolumn{1}{c}{Performance} & \multicolumn{1}{c}{BU-Net} &   \multicolumn{1}{c}{Fill holes}  &  \multicolumn{1}{c}{Random forest}  &  \multicolumn{1}{c}{Fill holes 2}\\
    \multicolumn{1}{c}{metric} &  \multicolumn{1}{c}{${\hat{\mathbf{S}}}^{BU}$} &   \multicolumn{1}{c}{$\hat{\mathbf{S}}^{H1}$} &  \multicolumn{1}{c}{$\hat{\mathbf{S}}^{RF}$} &  \multicolumn{1}{c}{$\hat{\mathbf{S}}^{H2}$} \\ \midrule
   Accuracy (\%) &  92.28 & 96.11 & 97.11 & 97.39\\
   Recall (\%) &  85.73 & 92.08 & 92.37 & 93.35 \\
   F1-score (\%) & 90.41 & 94.53 & 95.12 & 95.73 \\
   mIoU (\%) & 87.71 & 93.57 & 95.07 &  95.53 \\
\bottomrule
\end{tabular}}
\vspace{-0.25cm}
\end{table}

\begin{figure}[t!] %[htb]
\begin{minipage}[b]{1.0\linewidth}
  \centering
  \resizebox{8 cm}{!} {
  \centerline{\includegraphics[width=8.5cm]{Figures/Results/recall_boxplot.png}}}
  \vspace{-0.35cm}
\end{minipage}
\caption{\textbf{Box plot test results after each segmentation step.} }
\label{fig:boxplot}
\vspace{-0.4cm}
\end{figure}

% \begin{table}[t!]
% \centering
% \caption{\textbf{Test results after applying each step in our segmentation algorithm.} }
%     \vspace{-0cm}
%     \label{tab:seg-accuracy}
% \begin{tabular}{p{1.95cm}p{1.0cm}p{1.2cm}p{1.6cm}p{1.1cm}}
% %\begin{tabular}{ccccc}
% \toprule
%      \multicolumn{1}{p{1.95cm}}{Performance} & \multicolumn{1}{p{1.0cm}}{BU-Net} &   \multicolumn{1}{p{1.2cm}}{Fill holes}  &  \multicolumn{1}{p{1.6cm}}{Random forest}  &  \multicolumn{1}{p{1.1cm}}{Fill holes 2}\\
%     \multicolumn{1}{p{1.95cm}}{metric} &  \multicolumn{1}{p{1.0cm}}{${\hat{\mathbf{S}}}^{BU}$} &   \multicolumn{1}{p{1.2cm}}{$\hat{\mathbf{S}}^{H1}$} &  \multicolumn{1}{p{1.6cm}}{$\hat{\mathbf{S}}^{RF}$} &  \multicolumn{1}{p{1.1cm}}{$\hat{\mathbf{S}}^{H2}$} \\ \midrule
%    Accuracy (\%) &  92.28 & 96.11 & 97.11 & 97.39\\
%    Recall (\%) &  85.73 & 92.08 & 92.37 & 93.35 \\
%    F1-score (\%) & 90.41 & 94.53 & 95.12 & 95.73 \\
% \bottomrule
% \end{tabular}
% \vspace{-0.3cm}
% \end{table}

% \begin{table}[t!]
% \centering
% \caption{\textbf{Test results after applying each step in our segmentation algorithm.} }
%     \vspace{-0cm}
%     \label{tab:seg-accuracy}
% %\begin{tabular}{p{1.95cm}p{1.0cm}p{1.2cm}p{1.6cm}p{1.1cm}}
% \begin{tabular}{ccccc}
% \toprule
%      Performance & BU-Net &  Fill holes  &  Random forest  &  Fill holes 2\\
%     metric &  ${\hat{\mathbf{S}}}^{BU}$ &   $\hat{\mathbf{S}}^{H1}$ &  $\hat{\mathbf{S}}^{RF}$ & $\hat{\mathbf{S}}^{H2}$ \\ \midrule
%    Accuracy (\%) &  92.28 & 96.11 & 97.11 & 97.39\\
%    Recall (\%) &  85.73 & 92.08 & 92.37 & 93.35 \\
%    F1-score (\%) & 90.41 & 94.53 & 95.12 & 95.73 \\
% \bottomrule
% \end{tabular}
% \vspace{-0.3cm}
% \end{table}

%submitted
%variants 
%currently the 
%top performing method on the PASCAL
%report
%conducted 
%attaining
%focus
%yields
%significantly outperforming
%further improves performance to 63.1 percent. 
%observe any improvement when adopting either
%Employing the dense
%substantially outperforms the concurrent work 

% In Fig.~\ref{fig:seg-recall}, we analyze the performance dispersion of each step in our algorithm. We observe the skewness of the BU-Net performance: it is rather spread along the tail with lower performance, containing many instances with poorer accuracy and recall compared with their median. Note that once again the recall performance presents the worst results, with a lower bound extreme at 56.63\%. As we include the different steps in our algorithm, that dispersion is reduced notably. Therefore, despite not substantially improving the mean performance, the subsequent steps are the cornerstone to ensure the robustness of our algorithm, obtaining a final lower bound extreme of 88.57\% of recall.

% %Note that the whiskers are not actually at 1.5 x IQR from the box, they are the min and max data values within that range. 

% \begin{figure}[t!]
% \centering
% \includegraphics[width=3.45in]{Figures/Results/recall_boxplot.png}
%     \vspace{-0.3cm}
%     \caption{\textbf{Box plot test results after applying each segmentation step.} Performance dispersion of every algorithm step.}
%     \label{fig:seg-recall}
%     \vspace{-0.1cm}
% \end{figure}

\begin{table}[t!]
\centering
\caption{\textbf{Quantitative comparison on blade segmentation.} Relative metrics are with respect to the worst estimation.} 
\vspace{0cm}
\label{tab:seg-unet-compare}
\resizebox{8.5 cm}{!} {
\begin{tabular}{ccccccc}
\toprule
     \multicolumn{1}{c}{Method} & \multicolumn{1}{c}{Accuracy}  & \multicolumn{1}{c}{Recall} & \multicolumn{1}{c}{F1-score} & \multicolumn{1}{c}{mIoU} & Relative & Relative\\
    \multicolumn{1}{c}{} & {\%} & {\%}  & {\%}  & {\%} & accuracy & recall \\
    \midrule
   U-Net~\cite{unet} & 86.24 & 68.93 & 77.95 & 79.94 & 1 & 1\\ 
   DeepLabv3+~\cite{deeplabv3+}  & 94.14 & 87.38 & 89.03 & 90.74 & 1.09 & 1.27 \\ 
   SW~\cite{sw} & 93.48 & 91.71 & 91.37 & 89.71 & 1.08 & 1.33 \\ 
   ResNeSt~\cite{resnest} & 94.23 & 91.47 & 92.77 & 91.20 & 1.09 & 1.33 \\ 
   U-NetFormer~\cite{unetformer} & 96.20 & \textbf{93.51} & 94.42 & 93.65 & 1.12 & \textbf{1.36} \\ 
   Ours $\hat{\mathbf{S}}$ & \textbf{97.39} & 93.35 & \textbf{95.73} & \textbf{95.53} & \textbf{1.13} & 1.35 \\ 
\bottomrule
\end{tabular}}
\vspace{-0.5cm}
\end{table}



\textbf{Comparative Results.} As the BU-Net employs the U-Net architecture~\cite{unet}, we compared our proposed algorithm with the original one and state-of-the-art methods (Table~\ref{tab:seg-unet-compare}). We outperform those methods, thanks to overcoming the blade-background imbalances, and escaping from the tendency of predicting the most common region. 

%-----------------------------------------------------------------
\vspace{-0.3cm}
\subsection{Qualitative Results}
\vspace{-0.1cm}

We visualize the segmentation results of each algorithm step in Fig.~\ref{fig:seg-visual}. In general, our BU-Net can capture the blade edges and, therefore, its overall shape. However, it struggles with the inner area of the blade, which is addressed in the first hole filling step. Note that just using the standard hole filling approach would not improve the BU-Net outputs. Then, the random forest is responsible for fixing the misplaced pixels, acting as a regularizer of the pixel color intensity. By training a simple model, we ensure there is no overfitting and that it just learns the main patterns. Hence, if the BU-Net can in general classify which color intensities belong to the blade region, then the random forest would learn those patterns and resolve the particular images or pixels that failed, obtaining robustness in the solution. The second hole filling step is applied to solve the shortcomings the random forest can introduce. They are generally caused by dark regions that highly differ with the typical white color of the blade --a typical example is illustrated in Fig.~\ref{fig:seg-visual}, first row/right side--, where the random forest misplaces the vortex generators as background.

%the main patterns from all the images of the same blade surface, acting as a denoising step. In short, the random forest role is to tackle the specific errors of a single image. In the end, due to the simplicity of the random forest as a model

\begin{figure}[t!]
\resizebox{8.5 cm}{!} {
\hspace{-0.1cm}\begin{tabular}{@{}c|c@{}}
\includegraphics[width=0.4\textwidth]{Figures/Results/seg_visual_notruth_1.png}&
\includegraphics[width=0.4\textwidth]{Figures/Results/seg_visual_notruth_7.png}\\
%\includegraphics[width=0.4\textwidth]{Figures/Results/seg_visual_notruth_2.png}\\
\includegraphics[width=0.4\textwidth]{Figures/Results/seg_visual_notruth_3.png}&
\includegraphics[width=0.4\textwidth]{Figures/Results/seg_visual_notruth_4.png}\\
\includegraphics[width=0.4\textwidth]{Figures/Results/seg_visual_notruth_5.png}&
\includegraphics[width=0.4\textwidth]{Figures/Results/seg_visual_notruth_6.png}\\
%\includegraphics[width=0.4\textwidth]{Figures/Results/seg_visual_notruth_7.png}&
%\includegraphics[width=0.4\textwidth]{Figures/Results/seg_visual_notruth_8.png}
\end{tabular}}
\vspace{-0.2cm}
\caption{\textbf{Qualitative evaluation on test images after applying each step of our proposed segmentation algorithm.} On both sides, the same information is displayed. \textbf{First column:} Input color image ($\mathbf{I}$). \textbf{From second to fifth column:} BU-Net (${\hat{\mathbf{S}}}^{BU}$), first hole filling ($\hat{\mathbf{S}}^{H1}$), random forest ($\hat{\mathbf{S}}^{RF}$) and second hole filling ($\hat{\mathbf{S}}^{H2}\equiv\hat{\mathbf{S}}$) estimations, respectively.}
\label{fig:seg-visual}
\vspace{-0.35cm}
\end{figure}


In Fig.~\ref{fig:seg-fail}, we depict the most common cases where our segmentation algorithm could fail. To start with, the BU-Net might occasionally struggle with images of the wind turbine hub, because our dataset has a few instances of the hub. This is the case for the images in the first row. Distinct sources of contamination can also cause the BU-Net to mislead part of the blade with the background. This is the case for the image in the second row and right side. As the blade image borders are not fully captured and these images may contain a broad range of color intensities in the blade region, the hole filling and random forest steps are not sufficient to recover the blade region that has not been identified by the BU-Net.

In addition, an image could include wind turbines at its background or its own tower (see images in the second row). Ideally, we would like to predict the tower as blade, since it belongs to the turbine itself, but not the turbines in the landscape. This challenging problem is part of our future work. 


%Ideally, we would like to predict the tower as mask, since it belongs to the turbine itself. Turbines in the background of the image, however, should not be part of the blade. This challenging problem is part of our future work.



\begin{figure}[t!]
\resizebox{8.5 cm}{!} {
\hspace{-0.1cm}
\begin{tabular}{@{}c|c@{}}
\includegraphics[width=0.4\textwidth]{Figures/Results/seg_visual_fail_notruth_1.png}&
\includegraphics[width=0.4\textwidth]{Figures/Results/seg_visual_fail_notruth_2.png}\\
\includegraphics[width=0.4\textwidth]{Figures/Results/seg_visual_fail_notruth_3.png}&
\includegraphics[width=0.4\textwidth]{Figures/Results/seg_visual_fail_notruth_4.png}
\end{tabular}}
\vspace{-0.1cm}
\caption{\textbf{Failure cases on test images after applying each step of our segmentation algorithm.} See caption of Fig.~\ref{fig:seg-visual}.}
\label{fig:seg-fail}
\vspace{-0.0cm}
\end{figure}
%In both cases, the same information is displayed. \textbf{First column:} Input color image ($\mathbf{I}$). \textbf{From second to fifth column:} BU-Net (${\hat{\mathbf{S}}}^{BU}$), first hole filling ($\hat{\mathbf{S}}^{H1}$), random forest ($\hat{\mathbf{S}}^{RF}$) and second hole filling ($\hat{\mathbf{S}}^{H2}$) estimations, respectively.

%-----------------------------------------------------------
\vspace{-0.2cm}
\subsection{Windfarm Dissimilarity}
\vspace{-0.1cm}

%As mentioned in Section~\ref{sec:data}, as each dataset includes images from distinct windfarms, we examine in Fig.~\ref{fig:seg-inspections} how the performance varies depending on this factor.

%In each step of the algorithm, we observe significant dissimilar performance depending on the windfarm. For instance, in the first step when the BU-Net performance is reported, Windfarm \#1 performance is quite poorly compared with its piers, achieving approximately an accuracy of 83\%. In addition to that, each step of the algorithm can have drastically different effects between two windfarms. The extreme example comes again with Windfarm \#1, which after the first fill holing, it goes from the lowest accuracy to the highest performance. Conversely, for other cases such as Windfarm \#6, the first fill holing step does not signify any gain. A typical example from Windfarm \#1 that illustrates this phenomena is shown in Fig.~\ref{fig:seg-visual} (second row/ right side), where the BU-Net only captures the blade borders and then the hole filling is able to fill accurately the inner region of the blade. Another clear example comes with the random forest step, where while it significantly increases the accuracy of the Windfarm \#2, \#3 and \#7, it negatively affects the performance of Windfarm \#4 and \#5. An illustrative example of its improvement is shown in Fig.~\ref{fig:seg-visual} (first row/ left side).

Fig.~\ref{fig:seg-inspections} proves the generability of our approach, demonstrating that images from all the distinct windfarms have a high performance. Notice that this is not the case for the BU-Net without the post-processing steps, as it has really low performance for Windfarm \#1. Hence, it is the post-processing steps as a whole which ensures the generability. For instance, hole filling is crucial for Windfarm \#1, while random forest is crucial for Windfarm \#2 and \#3.
\vspace{-0.2cm}
%aquí resultados en el apéndice?

\begin{figure}[t!]
\centering
\includegraphics[width=3.1in]{Figures/Results/inspections.png}
    \vspace{-0.3cm}
    \caption{\textbf{Accuracy results over the test set for each windfarm.} It is included in each step in our method. For readability, the plot does not include every windfarm of the test set.}
    \label{fig:seg-inspections}
    \vspace{-0.3cm}
\end{figure}

%20210427-EUROWIND-DOSTRUP/Prise3200399T4-B-WW-pZR_001.jpg W7
%20210428-EUROWIND-DOSTRUP/Prise3200398T3-B-WW-tC3_010.jpg W2
%20211117-MDU-Diamond Hill/PriseTurbine07 Re-G-WW-DAw_010.jpg 
%Pioneer-Prairie-II/123 2S (119).JPG
%edf-turbine-2/B_WW_20155131003_backup.jpg
%edf-turbine-2/G_TE_20155131009_backup.jpg
%loegtved/a8ffa467-fd6e-49c8-9f1e-0d7f700bd72e_0044.jpg
%loegtved/a8ffa467-fd6e-49c8-9f1e-0d7f700bd72e_0078.jpg

%20210427-EUROWIND-DOSTRUP/Prise3200400T5-B-LE-XJr_001.jpg W7
%20211116-MDU-Diamond Hill/PriseTurbine12-R-LW-gFe_001.jpg
%Meridian-Way-II-207/207 2L (057).JPG W6
%upteko-ewii-marstrup/6_turbineSix_C_SS_0014.JPG








\vspace{-0.1cm}
\section{Conclusion}
\vspace{-0.1cm}

In this paper we have presented a novel segmentation algorithm for wind turbine blade images. To this end, we have proposed a BU-Net model in combination with additional steps based on hole filling and random forest strategies that can solve the problem in an accurate and efficient manner. Our BU-Net exploits a focal loss that is regularized with a contiguity term. The hole filling is simple but effective, and helps to improve the solution after applying every learning model. Finally, our unsupervised random forest is trained on-the-fly capturing the main patterns, like a denoising step. Remarkable quantitative and qualitative results are provided on newly acquired blade images, even when they belong to distinct inspection campaigns and wind farms, validating the capability of our model to generalize properly. Our future work is oriented to extend our model to semantic segmentation where different parts of the wind turbine need to be segmented in the image. Furthermore, we would like to study its impact on facilitating the defect detection task.

%Thanks to our segmentation algorithm, we believe our framework has the potential to be used for many applications for structural health monitoring. From a wind industry perspective, that means we can start to automate blade maintenance, saving costs and making green energy more competitive compared to traditional fuel power.

%This work exhibits the different strategies that were carried out to accomplish a novel image segmentation algorithm for wind turbine blade images. That is, a customized U-Net model complemented by two main post-processing approaches: an extended hole filling procedure and an unsupervised random forest trained on the fly. In particular, we regularized the focal loss with a contiguity term and we employed rotation-based data augmentation to define our BU-Net deep-learning model.

%Remarkable results are produced on new acquired blade images, even when they belong to distinct inspection campaigns and windfarms. This proves its capability to generalize properly. Furthermore, we believe our framework has the potential to be used for many applications for structural health monitoring, starting with facilitating the defect detection task. From a wind industry perspective, that means we can start to automate blade maintenance, saving costs and making green energy more competitive compared to traditional fuel power.






%compared to traditional fuel power
%submitted
%variants 
%currently the 
%top performing method on the PASCAL
%report
%conducted 
%attaining
%focus
%yields
%significantly outperforming
%further improves performance to 63.1 percent. 
%observe any improvement when adopting either
%Employing the dense
%substantially outperforms the concurrent work 
%examine
%cover
%introduce 
%measure
%suggest
%To this aim
%OPERATIONS MAINTENANCE





\vfill\pagebreak



% References should be produced using the bibtex program from suitable
% BiBTeX files (here: strings, refs, manuals). The IEEEbib.bst bibliography
% style file from IEEE produces unsorted bibliography list.
% -------------------------------------------------------------------------
\bibliographystyle{IEEEbib}
\bibliography{References}

\end{document}
